\documentclass[11pt,a4paper]{article}
\usepackage[margin=1in]{geometry}
\usepackage{amsmath,booktabs,hyperref,float}

\title{Iteration 015: Correlation Loss}
\author{Autoresearch Loop}
\date{April 8, 2026}

\begin{document}
\maketitle

\begin{abstract}
Training with Pearson correlation loss ($1-r$) instead of MSE achieves $r=0.373$,
not improving over the MSE-trained FIR ($r=0.376$). The model was trained on
correlation loss but early stopping used MSE validation, selecting a model from
the MSE-dominated warm-start phase rather than the correlation-optimized phase.
\end{abstract}

\section{Hypothesis}
Since the benchmark metric is Pearson $r$, directly optimizing correlation should
outperform MSE, which also penalizes scale/offset errors irrelevant to $r$.

\section{Method}
Loss: $\mathcal{L} = 1 - \bar{r}$, where $\bar{r}$ is mean Pearson correlation
across channels. Model: FIR (7 taps) + channel dropout. Early stopping on
validation MSE (this was the critical mistake).

\section{Results}
\begin{table}[H]
\centering
\begin{tabular}{lrrr}
\toprule
Model & Mean $r$ & Std $r$ & SNR (dB) \\
\midrule
FIR + dropout (iter011) & 0.376 & 0.076 & 0.61 \\
Correlation loss & 0.373 & 0.076 & 0.61 \\
\bottomrule
\end{tabular}
\end{table}

\section{Analysis}
The mismatch between training loss (correlation) and validation metric (MSE)
for early stopping was the key failure mode. Early stopping selected a model
from early training when MSE gradients still dominated, before the correlation
loss had time to refine the model. This motivates iter017: using correlation-based
validation for early stopping.

\section{Next}
Iterations 016--017: residual learning and combined loss with matched validation.

\end{document}
